\section{Определения}

\subsection{Сформулируйте теорему о приведении квадратичных форм к диагональному виду при помощи ортогональной замены координат.}
\begin{theorem}
    Любую квадратичную форму $q(x)$ можно ортогональным преобразованием привести к каноническому виду:
    \[
        q(y)=\lambda_1 y_1^2+\ldots+\lambda_n y_n^2,
    \]
    где $\lambda_i$ --- с.з. линейного оператора с той же матрицей, что и кв. ф..
    
    (Это называется приведением $q(x)$ к главным осям)
\end{theorem}
\subsection{Каков канонический вид ортогонального оператора? Сформулируйте теорему Эйлера.}

\begin{theorem}[О каноническом виде ортогонального оператора]
    Для любого ортогонального линейного оператора существует ОНБ, в котором матрица ортогонального оператора имеет блочно-диагональный (канонический) вид:
    \[
    A =
    \begin{pmatrix}
        A_{\varphi_1} &        &        &        &        &        & 0 \\
                      & A_{\varphi_2} &        &        &        &        &   \\
                      &        & \ddots &        &        &        &   \\
                      &        &        & A_{\varphi_k} &        &        &   \\
                      &        &        &        & 1      &        &   \\
                      &        &        &        &        & \ddots &   \\
        0             &        &        &        &        &        & -1
    \end{pmatrix},
\]
    где блоки имеют вид
    \[
        A_{\varphi_i}
        =
        \begin{pmatrix}
            \cos \varphi_i & -\sin \varphi_i \\
            \sin \varphi_i & \cos \varphi_i
        \end{pmatrix}.
    \]
    -- матрица поворота плоскости на угол $\varphi_i$.
\end{theorem}
\begin{theorem}[Теорема Эйлера]
    В $\mathbb{R}^3$ любое ортогональное преобразования обладает ОНБ, в котором его матрица имеет вид:
    \[
        A =
        \begin{pmatrix}
            \cos \varphi & -\sin \varphi & 0 \\
            \sin \varphi & \cos \varphi & 0 \\
            0 & 0 & \pm 1
        \end{pmatrix}.
    \]
\end{theorem}

\begin{remark}
    То есть любой ортогональный оператор в $\mathbb{R}^3$ является либо поворотом на угол $\varphi$
    (вокруг заданной оси-прямой, проходящей через $0$), либо композицией такого поворота и отражения.
\end{remark}


\subsection{Сформулируйте утверждение о QR-разложении.}

\begin{theorem}[О QR-разложении]
    Пусть
    \[
        A \in M_n(\mathbb{R})
    \]
    \[
        \det A \neq 0.
    \]
    
    Тогда существуют матрицы $Q$ и $R$ такие, что
    \[
        A = QR,
    \]
    где $Q$ --- ортогональная матрица
    \[
        Q^TQ = E,
    \]
    а $R$ --- верхнетреугольная матрица
    \[
        R =
        \begin{pmatrix}
            r_{11} & r_{12} & \ldots & r_{1n} \\
            0 & r_{22} & \ldots & r_{2n} \\
            \vdots & \vdots & \ddots & \vdots \\
            0 & 0 & \ldots & r_{nn}
        \end{pmatrix},
        \qquad
        r_{ii} > 0.
    \]
\end{theorem}

\subsection{Сформулируйте теорему о сингулярном разложении.}

\begin{theorem}[О сингулярном разложении]
    Для любой матрицы
    \[
        A \in M_{mn}(\mathbb{R})
    \]
    существует разложение
    \[
        A_{m \times n}
        =
        V_{m \times m}
        \Sigma_{m \times n}
        U^T_{n \times n},
    \]
    где
    \[
        U \in O_n(\mathbb{R}),
        \qquad
        V \in O_m(\mathbb{R}),
    \]
    то есть это ортогональные матрицы размеров $n \times n$ и $m \times m$ соответственно, а
    \[
        \Sigma
    \]
    --- диагональная матрица размера $m \times n$ с числами
    \[
        \sigma_1 \geq \sigma_2 \geq \ldots \geq \sigma_r
        >
        \sigma_{r+1}
        =
        \ldots
        =
        \sigma_{\min(m,n)} = 0,
    \]
    где
    \[
        r = RgA.
    \]
    
    Эти числа называются сингулярными числами матрицы $A$ и расположены по невозрастанию.
    
    То есть
    \[
        \Sigma =
        \begin{pmatrix}
            \sigma_1 &        &        &        & 0 \\
                     & \sigma_2 &        &        &   \\
                     &        & \ddots &        &   \\
                     &        &        & \sigma_r &   \\
            0        &        &  &        & \ddots   \\
                     &        &        &        &   &    0
        \end{pmatrix}_{m \times n}.
    \]
\end{theorem}

\subsection{Сформулируйте утверждение о полярном разложении.}

\begin{theorem}[О полярном разложении]
    Любой линейный оператор в евклидовом пространстве представим в виде композиции самосопряженного л.о. с неотрицат. с.з. и ортогонального л.о.,

    То есть для любой квадратной матрицы $A$ существует разложение
    \[
        A = S \cdot O,
    \]
    где $S$ --- симметрическая матрица с неотрицательными собственными значениями, а $O$ --- ортогональная матрица.
\end{theorem}
\subsection{Сформулируйте определение алгебры над полем. Приведите два примера.}

\begin{definition}
    Пусть $A$ --- векторное пространство над полем $\mathbb{F}$, снабженное операцией умножения векторов:
    \[
        *: A \times A \to A.
    \]
    
    $A$ называется алгеброй над полем $\mathbb{F}$, если выполнены следующие свойства:
    \[
        \forall x, y, z \in A,\quad \forall \alpha, \beta \in \mathbb{F}
    \]
    \[
        1)\ (x+y)*z = x*z + y*z
    \]
    \[
        2)\ x*(y+z) = x*y + x*z
    \]
    \[
        3)\ (\alpha x)*(\beta y) = (\alpha\beta)(x*y)
    \]
\end{definition}
\begin{example}
    $1) \quad
        \mathbb{C}
    $
    --- алгебра комплексных чисел над $\mathbb{R}$.
    
    При этом
    \[
        \text{Алгебра двумерная, }\dim\mathbb{C}=2
    \]

    $2) \quad
        M_n(\mathbb{F})
    $
    --- алгебра квадратных матриц над полем $\mathbb{F}$ с операцией умножения матриц.
    
    При этом
    \[
        \dim M_n(\mathbb{F})=n^2.
    \]
\end{example}
\subsection{Дайте определение эллипса как геометрического места точек. Выпишите его каноническое уравнение. Что такое эксцентриситет эллипса? В каких пределах он может меняться?}
\begin{definition}
    Эллипсом называют геометрическое место точек, сумма расстояний от которых до двух данных точек, называемых фокусами, постоянна.
\end{definition}

\begin{remark}
    Каноническое уравнение эллипса:
    \[
        \frac{x^2}{a^2} + \frac{y^2}{b^2} = 1.
    \]
\end{remark}
\begin{definition}
    Эксцентриситет эллипса:
    \[
        \varepsilon
        =
        \frac{\sqrt{a^2-b^2}}{a}
        =
        \sqrt{1-\frac{b^2}{a^2}}.
    \]
\end{definition}
\begin{remark}
    Эксцентриситет эллипса
    \[
        \varepsilon
        =
        \sqrt{1-\frac{b^2}{a^2}}
    \]
    лежит на полуинтервале
    \[
        [0,1)
    \]
    и служит мерой ``сплюснутости'' эллипса.
\end{remark}


\subsection{Дайте определение гиперболы как геометрического места точек. Выпишите её каноническое уравнение. Что такое эксцентриситет гиперболы? В каких пределах он может меняться?}
\begin{definition}
    Гиперболой называют геометрическое место точек, модуль разности расстояний от которых до двух данных точек, называемых фокусами, постоянен.
\end{definition}

\begin{remark}
    Каноническое уравнение гиперболы:
    \[
        \frac{x^2}{a^2} - \frac{y^2}{b^2} = 1.
    \]

\end{remark}
\begin{remark}
    Эксцентриситет гиперболы
    \[
        \varepsilon
        =
        \sqrt{1+\frac{b^2}{a^2}}
        >
        1
    \]
    (при $\varepsilon \to 1$ гипербола вырождается в два луча) характеризует угол между асимптотами.
\end{remark}
\subsection{Дайте определение параболы как геометрического места точек. Выпишите её каноническое уравнение. Что такое параметр параболы, каков его геометрический смысл?}
\begin{definition}
    Параболой называют геометрическое место точек плоскости, равноудаленных от данной точки (фокуса) и от данной прямой (директрисы).
\end{definition}
\begin{remark}
    Каноническое уравнение параболы:
    \[
        y^2 = 2px.
    \]
\end{remark}
\begin{remark}
    Число $p$, равное расстоянию от фокуса до директрисы (это его геометрический смысл), называют параметром параболы.
\end{remark}

\subsection{Сформулируйте теорему о классификации кривых второго порядка (нужны только названия возможных геометрических случаев).}

\begin{theorem}
    Для любой кривой второго порядка существует ПДСК $Oxy$, в которой уравнение этой кривой имеет один из следующих геометрических видов:
    \[
        1)\ \text{эллипс}
    \]
    \[
        2)\ \text{пустое множество}
    \]
    \[
        3)\ \text{точка}
    \]
    \[
        4)\ \text{гипербола}
    \]
    \[
        5)\ \text{пара пересекающихся прямых}
    \]
    \[
        6)\ \text{парабола}
    \]
    \[
        7)\ \text{пара параллельных прямых}
    \]
    \[
        8)\ \text{пустое множество}
    \]
    \[
        9)\ \text{прямая}
    \]
\end{theorem}

\subsection{Дайте определение цилиндрической поверхности. Приведите пример цилиндра второго порядка, отличного от эллиптического.}
Рассмотрим кривую $\gamma$, лежащую в некоторой плоскости $P$, и прямую $L$, не лежащую в $P$.

\begin{definition}
    Цилиндрической поверхностью называют множество всех прямых, параллельных $L$ и пересекающих $\gamma$.
\end{definition}

\begin{example}
    Гиперболический цилиндр:
    \[
        \frac{x^2}{a^2} - \frac{y^2}{b^2} = 1.
    \]
\end{example}

\begin{example}
    Параболический цилиндр:
    \[
        y^2 = 2px.
    \]
\end{example}

\subsection{Дайте определение линейчатой поверхности. Приведите два примера линейчатых поверхностей, не являющихся цилиндрическими поверхностями.}
\begin{definition}
    Линейчатой называют поверхность, образованную движением прямой линии.
\end{definition}

\begin{remark}
    Любой цилиндр является линейчатой поверхностью.
\end{remark}

\begin{example}
    Однополостный гиперболоид:
    \[
        \frac{x^2}{a^2}+\frac{y^2}{b^2}-\frac{z^2}{c^2}=1.
    \]
\end{example}

\begin{example}
    Гиперболический параболоид:
    \[
        \frac{x^2}{a^2}-\frac{y^2}{b^2}=2z.
    \]
\end{example}
\subsection{Запишите канонические уравнения эллиптического, гиперболического и параболического цилиндров. Для каждой поверхности указать на сколько частей она делит трехмерное пространство.}

\begin{remark}
    Эллиптический цилиндр:
    \[
        \frac{x^2}{a^2}+\frac{y^2}{b^2}=1
    \]
    делит трехмерное пространство на $2$ части.
\end{remark}

\begin{remark}
    Гиперболический цилиндр:
    \[
        \frac{x^2}{a^2}-\frac{y^2}{b^2}=1
    \]
    делит трехмерное пространство на $3$ части.
\end{remark}

\begin{remark}
    Параболический цилиндр:
    \[
        y^2=2px
    \]
    делит трехмерное пространство на $2$ части.
\end{remark}


\subsection{Запишите канонические уравнения эллипсоида, однополостного гиперболоида, двуполостного гиперболоида. Для каждой поверхности указать на какое число частей она делит трехмерное пространство.}

\begin{remark}
    Эллипсоид:
    \[
        \frac{x^2}{a^2}+\frac{y^2}{b^2}+\frac{z^2}{c^2}=1
    \]
    делит трехмерное пространство на $2$ части.
\end{remark}

\begin{remark}
    Однополостный гиперболоид:
    \[
        \frac{x^2}{a^2}+\frac{y^2}{b^2}-\frac{z^2}{c^2}=1
    \]
    делит трехмерное пространство на $2$ части.
\end{remark}

\begin{remark}
    Двуполостный гиперболоид:
    \[
        \frac{z^2}{c^2}-\frac{x^2}{a^2}-\frac{y^2}{b^2}=1
    \]
    делит трехмерное пространство на $3$ части.
\end{remark}


\subsection{Запишите канонические уравнения эллиптического параболоида, гиперболического параболоида. Для каждой поверхности указать на сколько частей она делит трехмерное пространство.}

\begin{remark}
    Эллиптический параболоид:
    \[
        \frac{x^2}{a^2}+\frac{y^2}{b^2}=2z
    \]
    делит трехмерное пространство на $2$ части.
\end{remark}

\begin{remark}
    Гиперболический параболоид:
    \[
        \frac{x^2}{a^2}-\frac{y^2}{b^2}=2z
    \]
    делит трехмерное пространство на $2$ части.
\end{remark}

\subsection{Дайте определение линейного функционала.}

\begin{definition}
    Отображение
    \[
        f: V \to \mathbb{F}
    \]
    называется линейной формой (линейным функционалом), если:
    \[
        1)\ f(x+y)=f(x)+f(y),\quad \forall x, y \in V
    \]
    \[
        2)\ f(\alpha x)=\alpha \cdot f(x),\quad \forall x \in V,\ \forall \alpha \in \mathbb{F}
    \]
\end{definition}

\subsection{Дайте определение сопряженного пространства.}
\begin{definition}
    Пространством, сопряженным к линейному пространству $V$, называется множество всех линейных форм на пространстве $V$ со следующими операциями:
    \[
        1)\ (f_1 + f_2)(x) = f_1(x) + f_2(x)
    \]
    \[
        2)\ (\alpha f)(x) = \alpha \cdot f(x)
    \]
\end{definition}

\begin{remark}
    Обозначение:
    \[
        V^* = \operatorname{Hom}(V, \mathbb{F})
    \]
    --- гомоморфизм линейных пространств
    \[
        V \to \mathbb{F}
    \]
\end{remark}

\subsection{Выпишите формулу для преобразования координат ковектора при переходе к другому базису.}

\begin{remark}
    Так как координаты линейной формы записывают в строку, то они преобразуются с помощью матрицы перехода
    \[
        T_{a \to b}, (\text{или } T_{a \to b}^T \text{ для транспонированного случая}),
    \]
    а не с помощью
    \[
        T_{a \to b}^{-1}.
    \]
    (то есть
    $
        [f]_b = [f]_a \cdot T_{a \to b}
    $
    или
    $
        [f]_b^T = T_{a \to b}^T [f]_a^T
    $
    Хотя для вектора $x \in V$ координаты меняются по формуле
    \[
        x^b = T_{a \to b}^{-1}x^a
    \]),
    то их (линейные формы) называются ковекторами, а "обычные"  векторы из $V$ --- контрвекторами.
\end{remark}

\subsection{Дайте определение взаимных базисов.}
\begin{definition}
    Базис
    \[
        e = (e_1,\ldots, e_n)
    \]
    в л.п. $V$ и базис
    \[
        f = (f^1,\ldots, f^n)
    \]
    в сопряженном пространстве $V^*$
    называются взаимными, если
    \[
        f^i(e_j)=\delta_j^i
        =
        \begin{cases}
            1,\quad i=j,\\
            0,\quad i\neq j.
        \end{cases}
    \]
\end{definition}